The intrinsic color indexes are $U - V = 0.300$
% sic: the source says "U - V = 0.300"; the problem gives B - V = 0.300
and $U - V = 0.330$, therefore
\[
U - B = (U - V) - (B - V) = 0.030
\]
The difference between the real and the measured color indexes are due to the
absorption in the interstellar medium, as well as due the stronger absorption
of the planetary nebula (in fact, during the calculations, one can easily see
that the interstellar absorption is not enough to account for the measured
values). So, for the first (nearest) star:
\begin{align*}
(B-V)_1 &= (B-V) + (E_B - E_V)\cdot D + (E'_B - E'_V)\cdot R = 0.327\\
(U-V)_1 &= (U-V) + (E_U - E_V)\cdot D + (E'_U - E'_V)\cdot R = 0.038\\
% sic: the values 0.038 and 0.365 on these two lines are swapped in the
% source; with the given data (U-V)_1 = 0.365 and (U-B)_1 = 0.038
(U-B)_1 &= (U-B) + (E_U - E_B)\cdot D + (E'_U - E'_B)\cdot R = 0.365
\tag*{(3 pt)}
\end{align*}
(the terms being, in order: apparent, intrinsic, interstellar absorption, and
nebula absorption). Calculating the interstellar absorption, we can find for,
the nebula values:
\begin{align*}
(E'_B - E'_V)\cdot R &= 0.0155\\
(E'_U - E'_V)\cdot R &= 0.0100\\
(E'_U - E'_B)\cdot R &= -0.0055 \tag*{(1 pt)}
\end{align*}
So, for the second (farthest) star,
\begin{align*}
(B-V)_2 &= (B-V) + (E_B - E_V)\cdot 3D + (E'_B - E'_V)\cdot 2R = \mathbf{0.3655}\\
(U-V)_2 &= (U-V) + (E_U - E_V)\cdot 3D + (E'_U - E'_V)\cdot 2R = \mathbf{0.425}\\
(U-B)_2 &= (U-B) + (E_U - E_B)\cdot 3D + (E'_U - E'_B)\cdot 2R = \mathbf{0.0595}
\tag*{(4 pt)}
\end{align*}
